\documentclass{article}
\usepackage[T1]{fontenc}
\usepackage[utf8]{inputenc}
\usepackage[polish]{babel}
\usepackage{url}
%{Informatyka stosowana 2020, I st., semestr VI}


\author{
	Mateusz Kosowski 251558 \\
	Nikodem Nowak 251598\\
dr inż. Marcin Kacprowicz
}

\title{Komputerowe systemy rozpoznawania 2024/2025\\Projekt 1. Klasyfikacja dokumentów tekstowych}
\begin{document}
\maketitle

Opis projektu ma formę artykułu naukowego lub raportu z zadania
badawczego (doświadczalnego/obliczeniowego0 (wg indywidualnych potrzeb związanych np. z
pracą inżynierską/naukową/zawodową). Kolejne sekcje muszą być numerowane i
zatytułowane. Wzory są numerowane, tablice są numerowane i podpisane nad
tablicą, rysunki są numerowane i podpisane pod rysunkiem. Podpis rysunku i
tabeli musi być wyczerpujący (nie ogólnikowy), aby czytelnik nie musiał sięgać do tekstu, aby go
zrozumieć.\\
\indent {\bf Limit stron sprawozdania: 20. Nie zmieniać parametrów tekstu
(czcionki, marginesów, interlinii, itp.). Istotne dla sprawozdania rysunki,
tabele, informacje, kody itp. można umieścić w załącznikach, które nie wliczają
się do limitu stron.}\\
\indent {\bf Wybrane sekcje (rozdziały sprawozdania) są uzupełniane wg wymagań w
opisie Projektu 1. i Harmonogramie Zajęć na WIKAMP KSR jako efekty zadań w~poszczególnych tygodniach}. 

\section{Cel projektu}
Badania będą prowadzone na gotowym zbiorze tekstów Reuters-21578 \cite{reuters}.
W celu przeprowadzenia obserwacji oraz wyciągnięcia wniosków zostanie napisana aplikacja. 
Jej zadaniem będzie ekstrakcja konkretnych cech tekstu, zamiana ich na wektory a następnie klasyfikacja w kategorii places, 
z możliwymi etykietami: west-germany, usa, france, uk, canada albo japan.


\section{Klasyfikacja nadzorowana metodą $k$-NN. Ekstrakcja cech, wektory cech}

\subsection{Zasada działania metody $k$-NN}
Algorytm $k$-NN (k najbliższych sąsiadów) \cite{knn} działa tak, że dla nowego, nieznanego dokumentu:
\begin{enumerate}
	\item Szuka w bazie $k$ tekstów najbardziej podobnych (najmniejsza odległość między wektorami cech).
	\item Sprawdza, do jakich kategorii należą znalezieni sąsiedzi.
	\item Przypisuje dokumentowi klasę najczęściej występującą wśród sąsiadów (głosowanie). W przypadku remisu stosuje odpowiednią metodę rozstrzygania.
\end{enumerate}

\subsection{Parametry wejściowe}
\begin{itemize}
	\item Wartość $k$ (liczba sąsiadów).
	\item Proporcje podziału tekstów na zbiór uczący i testowy (np. 60/40).
	\item Wybrane cechy oraz metryka odległości (np. euklidesowa).
\end{itemize}

\subsection{Wyniki i rezultaty}
Program dla każdego tekstu testowego zwraca przewidywaną etykietę (np. \texttt{usa}). Na podstawie porównania tych przewidywań z rzeczywistymi kategoriami, program oblicza ostateczne wyniki w postaci czterech miar jakości:
\begin{itemize}
	\item \textbf{Accuracy} (dokładność) – liczona dla całego zbioru dokumentów,
	\item \textbf{Precision} (precyzja) – liczona dla całego zbioru oraz oddzielnie dla każdej z 6 klas,
	\item \textbf{Recall} (czułość) – również dla całego zbioru i wybranych klas,
	\item \textbf{F1} – średnia harmoniczna miar Precision i Recall.
\end{itemize}
Dzięki tym parametrom można obiektywnie ocenić, jak dobrze model radzi sobie z przypisywaniem tekstów do odpowiednich państw.

\subsection{Wyodrębnione cechy (wektory cech)}
Każdy dokument reprezentowany jest przez wektor 10 cech:
\begin{enumerate}
	\item \textbf{Najczęściej występujące słowo} (Most Frequent Word) – wartość tekstowa, np. \texttt{trade}.
	\item \textbf{Najdłuższe słowo} (Longest Word) – wartość tekstowa, np. \texttt{internationalization}.
	\item \textbf{Średnia długość słowa} (Average Word Length) – wartość liczbowa, np. 5{,}43.
	\item \textbf{Bogactwo słownictwa} (Vocabulary Richness) – wartość liczbowa, np. 0{,}65.
	\item \textbf{Średnia długość zdania} (Average Sentence Length) – wartość liczbowa, np. 15{,}2.
	\item \textbf{Stosunek wielkich liter} (Uppercase Letter Ratio) – wartość liczbowa, np. 0{,}04.
	\item \textbf{Zagęszczenie znaków finansowych} (Financial/Numerical Sign Density) – wartość liczbowa, np. 0{,}012.
	\item \textbf{Indeks czytelności Flescha} (Flesch Reading Ease Index) – wartość liczbowa, np. 60{,}5.
	\item \textbf{Stosunek samogłosek do spółgłosek} (Vowel to Consonant Ratio) – wartość liczbowa, np. 0{,}73.
	\item \textbf{Suma wartości liczbowych} (Sum of All Numerical Values) – wartość liczbowa, np. 105{,}5.
\end{enumerate}

\subsection{Przykład wektora cech}
\begin{center}
$v = [$\texttt{trade}, \texttt{internationalization}, 5.43, 0.65, 15.2, 0.04, 0.012, 60.5, 0.73, 105.5$]$
\end{center}

\section{Miary jakości klasyfikacji} 
Miary jakości klasyfikacji (Accuracy, Precision,
Recall, F1). We wprowadzeniu zaprezentować minimum teorii potrzebnej do realizacji
zadania, tak by inżynier innej specjalności zrozumiał dalszy opis. Należy podać {\bf konkretne wzory miar użyte w tym eksperymencie oraz krótko
opisać ich znaczenie i zakresy przyjmowanych wartości. Należy podać przykładowe
wartości każdej miary. Nie przepisuj
teorii, ale podaj link/przypis i opisz jak rozumiesz jej zastosowanie w tym konkretnym
zadaniu}. \\
\indent Stosowane wzory, oznaczenia z objaśnieniami znaczenia symboli użytych w
doświadczeniu. Oznaczenia jednolite w obrębie całego sprawozdania.  Opis zawiera przypisy do bibliografii zgodnie z
Polską Normą, (zob. materiały BG PŁ).\\
\noindent {\bf Sekcja uzupełniona jako efekt zadania Tydzień 03 wg Harmonogramu Zajęć
na WIKAMP KSR.}


\section{Metryki i miary podobieństwa tekstów w klasyfikacji}
Wzory, znaczenia i opisy symboli zastosowanych metryk z
przykładami. Wzory, opisy i znaczenia miar
podobieństwa tekstów zastosowanych w obliczaniu metryk dla wektorów cech z
przykładami dla każdej miary \cite{niewiadomski08}.  Oznaczenia jednolite w obrębie całego sprawozdania.  {\bf Podaj metryki i miary
podobieństwa nie z literatury (te wystarczy zacytować linkiem), ale konkretne ich
postaci stosowane w zadaniu. Jakie zakresy wartości przyjmują te miary i
metryki, co oznaczają ich wartości? Podaj przykładowe wartości dla przykładowych wektorów cech}. \\ 
\noindent {\bf Sekcja uzupełniona jako efekt zadania Tydzień 04 wg Harmonogramu Zajęć
na WIKAMP KSR.}

\section{Wyniki klasyfikacji dla różnych parametrów wejściowych}
Wstępne wyniki miary Accuracy dla próbnych klasyfikacji na ograniczonym zbiorze tekstów (podać parametry i kryteria
wyboru wg punktów 3.-8. z opisu Projektu 1.). 
\noindent {\bf Sekcja uzupełniona jako efekt zadania Tydzień 05 wg Harmonogramu Zajęć
na WIKAMP KSR.}


\section{Dyskusja, wnioski, sprawozdanie końcowe}
We wnioskach z badania skomentować wyłącznie wpływ parametrów i cech na jakość procesu klasyfikacji. Innymi słowy – wnioski to to, co wynika z przeprowadzonych obliczeń, a nie to co wiadomo z literatury. 
Znaczenie kolejnych eksperymentów wg punktów 2.-8. opisu projektu 1.  Każdorazowo
podane parametry, dla których przeprowadzana eksperyment. 
Wykresy (np. słupkowe) i tabele wyników
obowiązkowe, dokładnie opisane w ,,captions'' (tytułach), konieczny opis osi i
jednostek wykresów oraz kolumn i wierszy tabel.\\ 

{**Ewentualne wyniki realizacji punktu 9. opisu Projektu 1., czyli ,,na ocenę 5.0'' i ich porównanie do wyników z
części obowiązkowej**.Dokładne interpretacje uzyskanych wyników w zależności od parametrów klasyfikacji
opisanych w punktach 3.-8 opisu Projektu 1. 
Omówić i wyjaśnić napotkane problemy (jeśli były). Każdy wniosek/problem powinien mieć poparcie
w przeprowadzonych eksperymentach (odwołania do konkretnych wyników: wykresów,
tabel). \\
\underline{Dla końcowej oceny jest to najważniejsza sekcja} sprawozdania, gdyż prezentuje poziom
zrozumienia rozwiązywanego problemu.\\

** Możliwości kontynuacji prac w obszarze systemów rozpoznawania, zwłaszcza w kontekście pracy inżynierskiej,
magisterskiej, naukowej, itp. **\\

\noindent {\bf Sekcja uzupełniona jako efekt zadań Tydzień 05 i Tydzień 06 wg Harmonogramu Zajęć
na WIKAMP KSR.}


\section{Braki w realizacji projektu 1.}
Wymienić wg opisu Projektu 1. wszystkie niezrealizowane obowiązkowe elementy projektu, ewentualnie
podać merytoryczne (ale nie czasowe) przyczyny tych braków. 

\begin{thebibliography}{0}

\bibitem{reuters}
\url{http://archive.ics.uci.edu/dataset/137/reuters+21578+text+categorization+collection}

\bibitem{knn}
Taunk, Kashvi; De, Sanjukta; Verma, Srishti; Swetapadma, Aleena. 
A Brief Review of Nearest Neighbor Algorithm for Learning and Classification. 
2019, pp. 1255--1260. DOI: 10.1109/ICCS45141.2019.9065747.

\end{thebibliography}

Literatura zawiera wyłącznie źródła recenzowane i/lub o potwierdzonej wiarygodności,
możliwe do weryfikacji i cytowane w sprawozdaniu. 
\end{document}
